%%
% 今時jarticleやjbook使ってる人いる？時代はjsarticleかjsbookだよ
% ついでに言うと、uplatexってのはplatexの上位互換、これを使わないなんて旧世代だよね
%
\documentclass[uplatex, report, a4j, 10pt]{jsbook}

\usepackage{packages/miyazaki-u-paper}   % 宮崎大学工学部の卒論の基本（片山先生作）を、僕がちょっと書き換えちゃった（テヘッ
\usepackage{enumitem}           % enumerate？古い古い
\usepackage[dvipdfmx]{graphicx} % 当然dvipdfmなんて使ってないよね
\usepackage[dvipdfmx]{color}    % listingsを使うときにはこれも必須、dvipdfmxを変えちゃうとgraphicxのdvipdfmxも変わるよ
\usepackage{listings, packages/jlisting} % コードを埋め込むなら必須
\usepackage{txfonts}            % フォントといえばやっぱりtxfonts、今はnewtxってのもあるらしい
\usepackage{verbatim}           % コメントアウトしてくれる便利なプリアンブルが使える \begin{comment} ... \end{comment}
\usepackage{url}
% \usepackage{easy-todo}
\usepackage[hdivide={21mm, , 21mm}, vdivide={30mm, , 25mm}]{geometry} % スタイルを少し変えたくても\hoffset, \voffsetは使わないでね
\usepackage{multirow}
\usepackage{ascmac}
\usepackage{subcaption}
\usepackage{colortbl}
\usepackage[dvipdfmx, hidelinks]{hyperref} % リンクを付けてくれる。
\usepackage{pxjahyper}          % リンクを付けてくれる（日本語）

% ソースコードの設定
\lstset{
  breaklines = true,%自動で折り返す
  basicstyle={\footnotesize\ttfamily},
  numberstyle={\scriptsize},
  stepnumber=1,
  numbersep=1zw,
  lineskip=-0.5ex,
  frame=single,
  numbers=left,%行番号を左に
  framexleftmargin=6mm,%行番号をフレーム内に
  numberstyle=\scriptsize,%行番号のサイズ
  stepnumber=1%1行おきに行番号を
}

% \addtolength\oddsidemargin{-1.5zw}  % 魔法の呪文01
% \addtolength\evensidemargin{-1.5zw} % 魔法の呪文02
% \addtolength\textwidth{3.0zw}       % 魔法の呪文03

\setcounter{page}{1}
\newcommand{\ttt}[1]{\texttt{#1}}
\newcommand{\ift}{if-then-else}
\newcommand{\tool}{SAMPLE\_APPLICATION}  % ツール名を設定

%%
% miyazaki-u-paper.sty用設定値
%
% \degree{g} % Graduateのg or Masterのm
% \figurenumbering{f} % 図目次を付ける場合はt (真) を持つ真偽値を引数に取る関数
% \tablenumbering{f} % 表目次を付ける場合はt (真) を持つ真偽値を引数に取る関数
\title{帳票画像の書類記入欄の自動検出による\\電子フォーム作成時間の削減}
\author{木村 優哉}
\nendo{4} % 年度
\advisor{片山 徹郎 教授} % 修論では無視する
\major{情報システム工学科}


\begin{document}
\maketitle

% \setlength\textfloatsep{0pt}% 魔法の呪文04

\preface{概要}
作成予定

%
% 本文
%


%%
% 謝辞
%

\section{はじめに}\label{cha:Introduction}
\tool 作ったよ。
先行研究\cite{example}を参考にしたよ。

\section{\tool の機能}

\section{\tool の適用例}
\tool にソースコード\ref{lst:input-example}を入力した結果を、図\ref{image/sample.png}に示す。
\begin{figure}[hb]
\begin{lstlisting}[label={lst:input-example}, caption={入力例}]
coordinate00,0,0
coordinate01,1,1
coordinate02,2,2
coordinate03,3,3
\end{lstlisting}
\end{figure}

\begin{figure}[hb]
    \centering
    \includegraphics[scale=0.5]{image/sample.png}
    \caption{\tool の画面}
    \label{image/sample.png}
\end{figure}

\section{\tool の評価}


\section{まとめ}

\footnotesize
%%
%参考文献

\begin{thebibliography}{0}
  \bibitem{example}Example Org: "Example Title". \url{http://example.com}. アクセス日: 2021/01/31.
\end{thebibliography}

\end{document}
